% Molt Petit — the from-scratch rewrite (2026-08).
% Ground rules for this document (agreed with George):
%  - written for the person who would DEPLOY the protocol; accurate, concise;
%  - target ~15 pages, no appendix unless a detail is truly needed but would
%    break the main flow;
%  - implementation + experiments get at most ONE page total, with one
%    concise measurements table; no slot-time formulas;
%  - every machine-checked claim cites its Lean name in the companion
%    codebase `Molt/` (which re-presents the model paper-first; theorems
%    are transported from the original development and axiom-guarded).
\documentclass[10pt]{article}

\usepackage[a4paper,margin=2.6cm]{geometry}
\usepackage{amsmath,amssymb,amsthm}
\usepackage{booktabs}
\usepackage{microtype}
\usepackage[hidelinks]{hyperref}
\usepackage{xcolor}
\usepackage{graphicx}

\definecolor{yakink}{HTML}{3A3733}
\definecolor{yakaccent}{HTML}{C3401F}

\newtheorem{theorem}{Theorem}
\newtheorem{assumption}{Assumption}
% \code: typewriter identifiers; long Lean names may break (hyphenless)
% after an underscore.
\DeclareRobustCommand{\code}[1]{{\renewcommand{\_}{\textunderscore\allowbreak}\texttt{#1}}}
\newcommand{\quorum}{\ensuremath{q}}
\newcommand{\fmax}{\ensuremath{f}}

\title{\textbf{Molt Petit: A Machine-Checked Slot-Based BFT Consensus
Protocol with Recursive Certificates}}
\author{George Agapov\\ \normalsize Yak Software, Andorra}
\date{September 2026, draft}

\begin{document}

\noindent{\LARGE\bfseries\raggedright\hyphenpenalty=10000 \exhyphenpenalty=10000
Molt Petit: A Machine-Checked Slot-Based BFT Consensus
Protocol with Recursive Certificates\par}
\vspace{4pt}
\noindent\begin{minipage}[b]{\dimexpr\textwidth-26mm\relax}
  George Agapov \textperiodcentered\ Yak Software, Andorra\par
  \vspace{2pt}
  {\small\color{yakink}September 2026, draft \textperiodcentered\ yak.finance/molt-petit\par}
\end{minipage}%
\hfill
\begin{minipage}[b]{20mm}
  \hfill\includegraphics[width=\linewidth]{yak-wordmark.pdf}%
\end{minipage}
\vspace{1em}

\begin{abstract}
Molt Petit is a slot-based Byzantine-fault-tolerant consensus protocol
whose entire light-client contract fits in one sentence: hold the genesis
and a clock, fetch a constant-size recursive certificate from anyone,
check that it verifies and that its tip is recent, and act on everything
$n$ blocks deep. Validity is a \emph{density} rule --- every matured
$n$-slot window must carry blocks from two thirds of the fixed roster ---
with no votes and deterministic, fixed-depth finality. All safety claims
are machine-checked in Lean~4 under the cryptographic and operational
assumptions of Section~\ref{sec:assumptions} --- the core claims about
the imported Rust validator source itself, proved sound against the model
and re-used as the recursive-proof circuit (with circuit-backend faithfulness
a stated trust assumption), and today's key-rotation results at
model level. Light-client safety is proved in a timed model against an
adversary that obtains signing keys at arbitrary real times --- by theft
or by transient compromise --- and signs whatever stamps they verify; the
fault budget counts only key exposure around each window's own time, so
the recency check is a machine-checked safety ingredient, and no valid
chain can run more than the budget's lookback ahead of real time --- and we develop key rotation inside
consensus: every mode recovers from lost keys, and against an adversary
who steals keys and keeps them forever, two modes come with exactly
stated operator duties and no client duty beyond the genesis and a clock;
for a public rotation schedule, theft safety is the light-client theorem
itself, with exposure read at the schedule. Measured with post-quantum signatures on commodity
hardware, the verified consensus rule adds under $2\%$ to any
certificate's proving cost.
\end{abstract}

\section{Introduction}
\label{sec:intro}

Consider a light client that is offline most of the time --- a wallet, an
embedded device, a bridge contract --- and occasionally needs to act on
the current state of a small consensus system. Today it either replays
history, trusts a service, or verifies a chain of committee hand-offs.
Molt Petit is a consensus protocol built so that this client can need
exactly three things: the deployment's genesis, a clock, and a
constant-size \emph{certificate} it can fetch from anyone. If the
certificate verifies and its tip is recent, the client may act on
everything $n$ blocks below it: any two such certificates agree at every
height at least $n$ below both tips, so two clients acting this way never
act on conflicting blocks --- and that guarantee is a machine-checked
theorem about the shipped validator code
(\code{Rust.rust\_timed\_certified\_agreement}; for the TypeScript
twin, \code{ts\_timed\_certified\_agreement}) under the assumptions of
Section~\ref{sec:assumptions}, not a design argument.

The protocol is deliberately small. A fixed roster of $n$ participants
produces blocks in public round-robin slots. There are no votes, no view
changes, and no fork-choice metric in the safety argument --- fork
choice is the trivial highest-valid-tip rule, load-bearing only for
liveness; the one structural rule is
\emph{density} --- a chain is valid only if every matured window of $n$
slots contains blocks from at least two thirds of the roster
(Section~\ref{sec:protocol}). Density makes deep forks impossible while
the fault budget holds, which turns finality into a fixed depth rather
than a probability. What a node holds is never history but a recursive
proof of the validity rule itself, plus a short suffix of signed blocks
(Section~\ref{sec:rules}).

Two ideas carry the security argument beyond the standard counting.
First, \emph{recency is a safety ingredient, not a freshness nicety}:
over a deployment's lifetime keys leak --- signers are compromised and
later cleaned, keys are stolen --- and a fork signed with keys exposed long
after its stamps is the classic long-range attack. The fault budget here
counts, for each window of $n$ slots, only the keys exposed during a
bounded stretch of real time around that window, from about $n$ slots
(plus the honest clocks' maximum lead) before it to $\varphi$ slots after it. A client that refuses tips older
than $\varphi$ is exactly the client for whom later exposure cannot
matter, and no adversary can prepare a valid chain further ahead of real
time than the lookback to get around the check (Section~\ref{sec:timed}). Second, \emph{key rotation lives inside
consensus}: every block declares its signing-key version, so rotation is
ordered by the chain like any other event. Every mode recovers from a lost key with no extra client duty.
Against an adversary who steals keys outright and keeps them forever,
two rotation modes --- a public schedule, and a lockstep counter with
erasure --- come with exactly stated operator duties, and their clients
may sleep arbitrarily long; the third, reactive rotation, recovers from
loss only, and we show why theft defeats it
(Section~\ref{sec:rotation}).

The engineering claim is as important as the theorems. The validator is
one Rust file; it is \emph{imported} into Lean~4 (via Charon + Aeneas)
and proved sound against the model, so the core safety theorems are
about the code that runs (the rotation pins are checked at model level
today, Section~\ref{sec:impl}); the same source is re-used as the recursive-proof
circuit, so the prover proves the predicate that was verified (circuit-backend
faithfulness being a stated trust assumption rather than a Lean-verified equivalence); and a
separate TypeScript implementation is proved sound against that
\emph{same} Lean model --- which is what keeps the model
implementation-agnostic while it is developed by agentic coding. A
build-enforced audit fixes what every headline theorem may assume: the
axioms of classical logic, and nothing else. Measured on
commodity hardware with post-quantum signatures, the verified consensus
rule is under $2\%$ of any certificate's proving cost
(Section~\ref{sec:impl}) --- formal consensus is effectively free on top
of the application.

The cost of this simplicity is stated up front and in
Section~\ref{sec:limitations}: a closed, equally weighted roster; no
liveness recovery (a starved window halts the chain by design); and
safety that depends on correct clocks rather than holding under
arbitrary network timing. Where those trades are
acceptable --- permissioned settings, bridges, embedded fleets ---
the result is a consensus layer whose entire client contract fits in one
sentence and whose proofs re-verify from a fresh checkout. The next
section states that contract completely, mode by mode, before any
definitions.

\section{How to run it: the operational model}
\label{sec:opmodel}

One protocol runs throughout this paper. A fixed roster of $n$
participants produces blocks in public round-robin slots; every block
declares the signing-key version it was signed under; a node holds a
constant-size certificate plus a short suffix of recent signed blocks;
and a verifier accepts a chain only if it validates and its tip is
recent --- at most $b$ slots old on the verifier's own clock --- then trusts everything $n$ blocks deep. Key rotation is part of the protocol
from the start: a deployment completes the validator by fixing one of
three rotation \emph{pins}. A pin bounds the key version a block may
declare --- a floor in modes 1--2, exact per window of a fixed grid in
mode 3 --- and so
decides when a retired key becomes dead; that choice is the mode (1--3
below). Mechanically the modes
differ only in what computes the pin; in what they guarantee, they
differ only against key \emph{theft}. Producers rotate; clients never
react to rotations --- each client duty below is a fixed rule in time,
never a response to an event. (Names in \code{typewriter} are
Lean theorems and definitions; every one resolves in the companion
library \code{Molt/}, except the per-implementation bridge names of
Section~\ref{sec:impl}.)

\paragraph{Without theft, every mode is Theorem~1.}
Every adversary in this paper can make a producer \emph{lose} its key
--- destroy it, wedge its signer, force a rotation. Against an adversary
with only that power, all three modes give the same guarantee to the
same client: \textbf{hold the genesis and a clock; sleep as long as you
like}. Loss changes none of the hypotheses of Theorem~\ref{thm:lc}, and
each mode's validator inherits it directly (\code{keyrot\_loss\_agreement},
\code{sched\_loss\_agreement}, \code{lockstep\_loss\_agreement};
Theorem~\ref{thm:keyloss}). The loser recovers by its mode's own
discipline --- an in-band rotation in mode 1, a later scheduled
generation in mode 2, the next lockstep switch in mode 3 --- at the
price of its missed slots (mode 1's local single-block production step
under continued density coverage is checked in the core,
\code{liveness\_produce\_blockK}). The modes differ only against key
\emph{theft}, which is what the rest of this section is about.

\paragraph{Mode 1: reactive rotation --- recovery from key loss, not theft.}
Producers rotate whenever they choose, routinely or after losing a key,
and \emph{no one knows in advance when rotations will unfold} --- there
is no schedule, public or private. The client holds the genesis and a
clock and gets exactly the no-theft guarantee above
(Theorem~\ref{thm:keyloss}), and nothing more. Mode 1 is \textbf{not safe against key
theft}: an adversary that steals the retired keys of a quorum of seats
can forge a recent fork that branches before their rotations
(Section~\ref{sec:rotation}). A deployment that must survive theft runs
mode 2 or mode 3.

\paragraph{Mode 2: scheduled rotation --- for clients that sleep.} The
key version in force at any slot is a \emph{public function of the
slot}, fixed and published before the deployment starts: every producer,
verifier, and client knows the entire rotation schedule in advance, and
the validator refuses versions below it. Every mode signs under keys
delegated from one offline master key (a \emph{cold root}) per seat
(Section~\ref{sec:protocol}); what mode~2 adds is \emph{when}: the
operator derives each period's signing key --- each \emph{generation}
--- just before its scheduled span (its \emph{era}) begins, a
discipline the theft theorem charges through its exposure budget rather
than assuming (Section~\ref{sec:rotation}). The client holds the genesis, a clock,
and the constant that defines the schedule; it may stay offline
arbitrarily long and sync whenever it wakes. No anchor, no cadence, no erasure.

\paragraph{Mode 3: free-cadence lockstep --- for clients that sleep and
operators that erase.} No schedule anywhere, not even a constant. The
roster advances one shared generation counter at moments of its own
choosing; validity forbids mixing generations within any window of the
fixed window grid --- slots $[Wn,\,Wn{+}n)$ --- and never lets the
counter decrease across grid windows; and retired key material is
physically destroyed at each switch (hardware modules or RAM-only keys
--- storage that does not truly erase does not qualify). The client
holds the genesis and a clock --- even less than mode 2 --- and may
sleep arbitrarily long. The price sits entirely with the operators:
lockstep coordination, and erasure that is real. (Mode 3's guarantees hold at the certificate presentation as well as for
full chains, under the cumulative budget at depth $n$
(\code{lockstep\_recent\_certified\_suffix\_agreement}) and under the
per-generation census at depth $2n$
(\code{lockstepGen\_recent\_certified\_suffix\_agreement}).
Erasure's per-generation credit against the budget is machine-checked:
under a per-generation census --- for each generation $j$ separately, at
most $T$ seats whose generation-$j$ key is stolen --- the same agreement
holds at confirmation depth $2n$ instead of $n$ (at equal tip heights,
\code{lockstep\_client\_safety\_gen}; at unequal heights the lower
chain's $2n$-deep block lies on the taller,
\code{lockstepGen\_recent\_tip\_ancestor\_mem}), and with no
shared-genesis hypothesis; the
cumulative, depth-$n$ form remains available. Section~\ref{sec:rotation}.)

\paragraph{The modes at a glance.}
\begin{center}
\small
\setlength{\tabcolsep}{3.5pt}
\begin{tabular}{@{}p{1.4cm}p{2.1cm}p{2.3cm}p{1.5cm}p{1.0cm}p{5.8cm}@{}}
\toprule
 & pin computed from & client holds & erasure & anchor & headline Lean theorem \\
\midrule
no theft, any mode & any of the three &
  genesis + clock & no & no &
  Thm.~\ref{thm:lc}: \code{timed\_light\_client\_safety}
  (+~Rust: \code{Rust.rust\_timed\_certified\_agreement}); per validator
  \code{keyrot\_loss\_agreement}, \code{sched\_loss\_agreement},
  \code{lockstep\_loss\_agreement} \\
\addlinespace
1 reactive & chain floor, $n$ slots deep &
  --- & --- & --- &
  none against theft (Section~\ref{sec:rotation}); loss: row above \\
\addlinespace
2 scheduled & slot arithmetic $\mathit{gen}(s)$ &
  genesis + clock + $\mathit{gen}$ & no & no &
  Thm.~\ref{thm:sched}: \code{sched\_exposure\_agreement},
  \code{sched\_exposure\_certified\_agreement} \\
\addlinespace
3 lockstep & roster counter, no-mixing &
  genesis + clock & yes (credited, Thm.~\ref{thm:lock-gen}) & no &
  Thm.~\ref{thm:lock}: \code{lockstep\_client\_safety};
  Thm.~\ref{thm:lock-gen}: \code{lockstep\_client\_safety\_gen},
  \code{lockstepGen\_recent\_tip\_ancestor\_mem} \\
\bottomrule
\end{tabular}
\end{center}
Read the trust downward: mode 1 trusts that keys are lost, never
stolen; mode 2 trusts slot arithmetic, with key-provisioning discipline
charged to its budget;
mode 3 trusts roster coordination plus physical destruction. Mode 2
runs on Theorem~\ref{thm:lc}'s own windowed budget, read at the
schedule; mode 3 on the $\rho + T \le \fmax$ budget --- its $T$ may be
read per generation, at depth $2n$ --- and a recency-scoped signature
surface; every mode without theft is the no-theft row.

Everything that follows makes these paragraphs precise: the protocol and
its validator (Sections~\ref{sec:protocol}--\ref{sec:rules}), what a
deployment must assume (Section~\ref{sec:assumptions}), and the theorems
each contract rests on (Section~\ref{sec:results}).

\section{The protocol}
\label{sec:protocol}

Machine-checked claims cite their Lean names in \code{typewriter}; every
cited name resolves in the companion library \code{Molt/}, except the
per-implementation bridge names of Section~\ref{sec:impl}, which live
with their pipelines.

\subsection{Setting}

Time is divided into slots of a fixed duration $\tau$. A fixed, closed set
of $n$ participants --- the \emph{roster} --- produces blocks. Slot $s$
belongs to roster seat $s \bmod n$ (\code{producer}): who may produce, and
when, is deterministic and public; there is no leader election. Two
constants derived from $n$ appear throughout: the \emph{quorum}
$\quorum = \lceil 2n/3 \rceil$ and the \emph{fault budget}
$\fmax = \lfloor (n-1)/3 \rfloor$ --- the number of slots per $n$-slot
window whose keys the adversary may hold
(Assumption~\ref{ass:budget} makes this precise).
Table~\ref{tab:notation} collects the paper's symbols.

\begin{table}[t]
\centering
\small
\begin{tabular}{@{}p{3.4cm}p{9.6cm}l@{}}
\toprule
symbol & meaning & where \\
\midrule
$n$ & roster size; also the density window (slots) and trust depth
  (blocks) & \S\ref{sec:protocol} \\
$\quorum = \lceil 2n/3 \rceil$ & window quorum & \S\ref{sec:protocol} \\
$\fmax = \lfloor (n{-}1)/3 \rfloor$ & per-window fault budget & Asm.~\ref{ass:budget} \\
$\tau$ & slot duration (seconds) & \S\ref{sec:impl} \\
$\mathit{now}$ & verifier's clock; recency test
  $\mathit{now} \le \mathit{tip.slot} + b$ & \S\ref{sec:protocol} \\
$\varphi$, $\ell$, $\sigma$ & freshness $b + \delta$; budget lookback;
  honest clock lead & Asm.~\ref{ass:budget}--\ref{ass:clock} \\
$H(\cdot)$ & the block-id hash & \S\ref{sec:protocol} \\
$\mathit{keyFor}(i,j)$, \code{keyIndex} & seat $i$'s public key at
  version $j$; the version a block declares & \S\ref{sec:protocol} \\
$\rho$, $T$ & per-window rented slots and stolen-key producers,
  $\rho + T \le \fmax$ (mode 3) & \S\ref{sec:results} \\
$E$, $\mathit{gen}(s) = \lfloor s/E \rfloor$ & rotation-schedule era length;
  scheduled key version (mode 2) & \S\ref{sec:results} \\
\bottomrule
\end{tabular}
\caption{Notation. The density window, the trust depth, and the
confirmation depth at which a mode-1 rotation takes force
(Section~\ref{sec:results}) are all set by $n$. The recency bound $b$ is
the deployment's choice; it sets how long after a window key exposure
still counts against it (Assumption~\ref{ass:budget}).}
\label{tab:notation}
\end{table}

\subsection{Blocks and chains}

A block is six fields:
\[
B = (\mathit{slot},\ \mathit{height},\ \mathit{prev},\ \mathit{id},\
\mathit{contentsHash},\ \mathit{keyIndex}).
\]
The payload (transactions, in a deployment) travels out of band;
\code{contentsHash} commits to it, and no consensus rule ever looks inside
it. \code{keyIndex} is the version of the signing key the producer used:
key rotation is part of the consensus state, carried in band, and
Section~\ref{sec:results} builds on it. Every block travels with its
producer's signature over it --- carried alongside the six fields, never
inside them. \code{prev} is the parent's id, and \code{id} is a hash of
the block itself, computed and checked by the deployment's hashing layer
under a fixed contract:
\[
\mathit{id} \;=\; H(\mathit{slot},\ \mathit{height},\ \mathit{prev},\
\mathit{parentSig},\ \mathit{contentsHash},\ \mathit{keyIndex}).
\]
The consensus rules treat ids as opaque. Two preimage choices do real work
later. The \emph{parent's signature} is inside the hash: that is what
forces anyone fabricating a chain to produce its signatures one after
another rather than all at once
(Section~\ref{sec:results} turns this into a wall-clock bound).
And the \emph{key version} is inside the hash, so two blocks differing
only in it hash differently --- without that, the hash assumption becomes
unsatisfiable the moment any producer has rotated.

A chain is a sequence of blocks with heights increasing by one, slots
strictly increasing, and each block linking to its parent's id, rooted
at a genesis block --- height zero, no parent (\code{linksOk},
\code{genesisOk}); that the root is the deployment's genesis is what
the theorems hypothesize (Section~\ref{sec:assumptions}).

\subsection{The density rule}

There are no votes and no longest-chain comparison. A chain must be
\emph{dense} (\code{validChain} names the combined structure-and-density
checks): every matured window of $n$
consecutive slots must contain at least $\quorum$ blocks of the chain,
where a window $[u, u+n)$ has matured once the tip has reached its end,
$u + n \le \mathit{tip.slot} + 1$ (\code{denseSoFar}). Read it
operationally: a chain is acceptable evidence only if, at every point of
its history, at least two thirds of the roster was actively extending it.

Density is the whole safety mechanism. If two chains diverge and both stay
dense, each must collect $\quorum$ blocks in any window past the
divergence; since $2\quorum > n + \fmax$, some slot in that window carries
a block of \emph{both} chains and belongs to an honest producer --- and an
honest producer extends one chain per slot, never two. This counting
argument is the core combinatorial engine of every safety theorem in
Section~\ref{sec:results}, around which cryptographic unforgeability,
certificate grounding, and rotation pin mechanisms are layered and
machine-checked.

Chain validity has one more condition, riding on the \code{keyIndex}
field: along the chain, a producer's declared key version never
decreases (\code{keyMonoOk}) --- rotation only ever moves forward, and
the chain's order is the agreed order in which rotations happen.
Structure, density, and this monotone rule together are
\code{validChainK}: \emph{the} chain validity of this paper. There is
exactly one protocol here, with key rotation built in from the start; a
deployment completes it by fixing one of three rotation \emph{pins}
(Section~\ref{sec:rotation}).

\subsection{What a node holds}

A node does not hold history. It holds a \emph{certified chain}: a
constant-size certificate plus a short suffix of recent signed blocks. The
certificate attests everything below the suffix and exposes a
\emph{claim}: the tip slot, height, and id it vouches for, plus a small
boundary buffer of near-tip blocks so that window density can be counted
across the seam. In the flagship implementation the certificate is a
\emph{recursive} zero-knowledge proof: each new proof checks the
previous proof plus the newly folded blocks, so one constant-size object
ends up attesting arbitrarily long history.

The rules touch cryptography only through two small interfaces that a
deployment implements:
\begin{itemize}
\item \code{SigOps} --- $\mathit{sign}$ and $\mathit{verify}$ ---
  together with the versioned key map $\mathit{keyFor}(i,j)$, the
  public key of seat $i$'s signing key at version $j$ (a separate
  parameter, \code{KeyRegistry}, in the model; the implementation's
  dictionary folds it into \code{SigOps});
\item \code{CertOps} --- verify a certificate, read its claim, and
  extend it by one validated block; the implementation's dictionary
  (Section~\ref{sec:impl}) also exposes a by-tip-id lookup of locally
  held certificates, which compaction uses.
\end{itemize}
One deployment note on $\mathit{keyFor}$, because it is where a trusted
directory service would creep in: there is none, and a verifier looks
nothing up. $\mathit{keyFor}$ is a mathematical parameter of the model
--- which public key each $(\text{seat},\text{version})$ pair denotes
--- and a block carries its own witness for it. A block's
``signature'' is in practice a bundle --- the signing public key, a
delegation certificate binding that key to seat $i$ \emph{and version
$j$}, and the signature proper --- verified against a fixed set of
\emph{cold root} keys, one per roster seat, created at the deployment
ceremony. The delegation certificate must bind the version: one delegate
key answering for several versions is exactly the key reuse that
Section~\ref{sec:results} has to exclude by assumption.

How the theorems read $\mathit{verify}$. The model's signature predicates
(``$B$ carries a verifying signature'') are existential over signatures,
so they are read against the ideal signature functionality
$\mathcal{F}_{\mathrm{SIG}}$~\cite{canettisig} at the verifier's real slot: a
signature verifies exactly when it has been produced by then --- by an
honest signer, or by whoever holds an exposed key. A scheme realizes this
reading precisely when it is EUF-CMA~\cite{canettisig}. The reading matters
for one hypothesis: \code{hbridge} (``every admissible block was produced
by real slot $R$'') is then plain causality, whereas against a concrete
scheme's $\mathit{verify}$, which accepts signatures nobody has computed
yet, it would be false.

\subsection{The verifier rule}
\label{sec:verifier-rule}

Everything a consumer of chains --- above all, a light client --- does is
one rule. Accept a certified chain if and only if:
\begin{enumerate}
\item it \emph{validates}: the five checks of Section~\ref{sec:rules}
  --- certificate, signatures at declared versions, linkage, density,
  and the key-version discipline with the deployment's pin
  (\code{validCertifiedChain} plus the mode's key checks);
\item it is \emph{recent}: $\mathit{now} \le \mathit{tip.slot} + b$ on the
  verifier's own clock, where the real time $R$ is at most $\delta$ slots ahead of
  $\mathit{now}$ ($R \le \mathit{now} + \delta$); so the tip is at most
  $\varphi = b + \delta$ slots old in real time.
  (This clock comparison is the one line a
  deployment implements outside the verified validator ---
  Section~\ref{sec:limitations}.)
\end{enumerate}
It then trusts the blocks $n$ deep below the tip. Both halves are
necessary: validity without recency admits forks signed with keys
exposed long after the fact --- the long-range attack
(Section~\ref{sec:results}) --- and recency without validity means
nothing. The parameter $n$ governs the density window and the trust
depth. The recency bound $b$ is traded against the fault budget: a tip
may be up to $\varphi = b + \delta$ slots old in real time, and key
exposure counts against a window until $\varphi$ slots after it
(Assumption~\ref{ass:budget}). For example, $b = n$ at the prototype's
parameters ($n = 7$, $\tau \approx 4$\,s) accepts tips up to about half a
minute old on the verifier's clock.

\section{The rules a node runs}
\label{sec:rules}

A node's whole life is three operations on its certified chain. All three
are pure functions --- deterministic, no I/O, no hidden state --- which is
what lets one source of truth serve as reference implementation, Lean
model, and circuit statement at once (Section~\ref{sec:impl}).

\paragraph{Validate} (\code{validCertifiedChain}). Accept a certified
chain only if all of the following hold:
\begin{enumerate}
\item the certificate verifies (\code{CertOps.verify} --- one cheap check,
  e.g.\ of a recursive proof);
\item every suffix block's signature verifies under
  $\mathit{keyFor}(\text{slot's producer},\ \text{declared version})$;
\item the first suffix block extends the claim's tip: height one up, slot
  strictly later, \code{prev} equal to the tip id --- and every adjacent
  suffix pair links the same way;
\item every window that matures within the suffix is dense, counted over
  the claim's boundary buffer plus the suffix (\code{validSuffix});
\item the key-version discipline holds: a producer's declared versions
  never decrease (\code{keyMonoOk}, Section~\ref{sec:protocol}), and
  every block's declared version satisfies the deployment's rotation
  \emph{pin} (a floor in modes 1--2, exact per grid window in mode 3) --- the
  one mode-specific check, fixed at deployment. These two are
  stated over the full-chain form of the validator
  (\code{validSignedChainK$'$}; the same validator under the mode 2--3
  pins); at certificate level they are proved for modes 2 and 3
  (\code{sched\_exposure\_certified\_agreement},
  \code{lockstepGen\_recent\_certified\_suffix\_agreement}).
\end{enumerate}
The model replays the structural checks from genesis
(\code{validChain}); a deployed
node checks only the windows that mature with each arriving block. The two
are equivalent --- earlier windows were checked when the parent was
validated --- so the incremental check is an optimization, not a
loosening.

\paragraph{Produce} (\code{produceBlock?}). Only in the caller's own slot:
build the next block on the current tip, sign it, append it to the suffix
--- and ship the result only if it passes the same \code{validSuffix}
every other node runs. That final gate is the point: nothing on the
production path has to be trusted for safety, because anything malformed
it could emit is exactly what every validator rejects.

\paragraph{Select} (\code{selectChain}). Adopt a candidate chain only if
it validates and its tip is strictly higher than the current one's.
Selection never carries safety: the theorems of
Section~\ref{sec:results} hold for \emph{any} two chains a verifier
accepts, however they were selected. Selection matters only for
liveness.

\paragraph{Compaction.} Periodically a node folds the oldest suffix
blocks into the certificate and keeps the suffix short. Certificates
carry no authority --- they are proofs --- so any node may run the
prover, and one untrusted proving service suffices; if proving lags, the
suffixes simply grow until it catches up, and validity is unaffected.
Because certificates arrive in batches, compaction scans the possible cut
points from newest to oldest and jumps at the first position for which it
holds a verified certificate whose claim equals the current claim
advanced by exactly the dropped blocks. The compacted chain must then
pass \code{validCertifiedChain} like anything else; compaction, like
production, is outside the trusted base.

\section{What a deployment must provide}
\label{sec:assumptions}

Every assumption below except the clock appears literally as a
hypothesis of the theorem that cites it --- Assumptions~\ref{ass:budget}--\ref{ass:cert}
in Theorems~\ref{thm:lc}--\ref{thm:forge}, Assumption~\ref{ass:delivery}
in the liveness theorem. Mode 2's theft theorem keeps them, reading
exposure at the schedule's key versions and restricting
Assumption~\ref{ass:sig} to the blocks its validator admits; mode 3
replaces Assumptions~\ref{ass:budget}--\ref{ass:sig} with rescoped forms
defined in Section~\ref{sec:rotation}'s adversary paragraph. Lean's axiom audit
confirms that nothing further is assumed behind the scenes. The clock is the one
exception by nature: Theorem~\ref{thm:lc} is stated at the real slot $R$
at which the verifier checks, the rotation theorems at the verifier's
reading $\mathit{now}$, and Assumption~\ref{ass:clock} says how the two
are related.

The adversary, stated once: message delivery is entirely adversarial for
every safety result (timeliness enters only for liveness,
Assumption~\ref{ass:delivery}); the adversary obtains signing keys --- by
stealing them, or by controlling a signer for a while --- at any real
times it likes, within the budget below, and an exposed key signs
anything, stamped with any slot, at any real time; the verifier's clock
satisfies the real-time bound $R \le \mathit{now} + \delta$
(Assumption~\ref{ass:clock}) while the adversary's timing is
unconstrained. What the adversary cannot do is exactly
Assumptions~\ref{ass:sig}--\ref{ass:cert}: sign under keys that are not
exposed, find hash collisions among blocks that occur, or forge
certificates.

\begin{assumption}[Exposure budget]\label{ass:budget}
The key that verifies stamp $s$ (seat $s \bmod n$'s) is \emph{exposed} at
real slot $r$ when someone other than its honest holder can sign under it
then: the seat is Byzantine, the key has been stolen, or the adversary
controls the signer. Exposure is arbitrary in time --- a stolen key stays
exposed, control of a signer may come and go. The budget: of any $n$
consecutive stamps $[v, v + n)$, at most $\fmax = \lfloor (n-1)/3
\rfloor$ have their key exposed at some real slot in
$[v - \ell,\ v + n + \varphi)$, where $\varphi$ is the freshness of
Assumption~\ref{ass:clock} and the lookback $\ell$ satisfies
$\ell \ge n - 1$ and $\ell \ge n + \fmax + \sigma + 1 - \quorum$ for the
honest clock lead $\sigma$ of Assumption~\ref{ass:sig}. For
$n = 3\fmax + 1 \ge 4$, $\ell = n - 1 + \sigma$ suffices: the stretch runs
from $n - 1 + \sigma$ slots before the window to $\varphi$ slots after it. Exposure outside that stretch of real
time --- control that ended before it, theft that happened after it ---
does not count against the window.
\quad\textit{(Lean: \code{ExposureBounded}.)}
\end{assumption}

\begin{assumption}[Signatures]\label{ass:sig}
Signing happens at real slots, and the assumption speaks only of
\emph{admissible} blocks: those the verifier's per-block check accepts ---
the producer signature verifies, and the id is \emph{formed}, that is, it
is the hash of the block's preimage taken over its parent's actual
signature, as the deployment's hashing layer recomputes it
(Section~\ref{sec:protocol}) (\code{SigningExecutionOn}). The assumption
bundles: \emph{(a) custody} --- while a stamp's key is unexposed, only its
honest holder signs under it; \emph{(b) unforgeability} --- no one can
produce a valid signature on a message the key holder never signed, even
after obtaining signatures on messages of their own choosing (the
standard EUF-CMA notion); and \emph{(c) the honest signing discipline}
--- an honest holder signs at most one block per stamp, ever, keeping a
persisted signing record as the standard guard against equivocation
after a restart (with (a): \code{honest\_once}); its clock leads real time
by at most $\sigma$ slots, so it never signs a stamp more than $\sigma$
slots ahead (\code{HonestClockOn}); and every admissible block accepted
by the verifier at real slot $R$ was produced at some real slot no later
than $R$ (\code{hbridge}, causality plus EUF-CMA, under the reading of
$\mathit{verify}$ in Section~\ref{sec:protocol}).
The id-formation contract of Section~\ref{sec:protocol} ensures that an
admissible block is signed only once its parent is available
(\code{chain\_order}): its signature covers its id, whose preimage
contains the parent's signature.
\quad\textit{(Lean: \code{SigningExecutionOn}, \code{HonestClockOn}, \code{hbridge}.)}
\end{assumption}

Nothing constrains an exposed key: it signs blocks stamped for past or
future slots, at any real time, in any number; the only limit is how many
keys are exposed, and when (Assumption~\ref{ass:budget}). A block it signs
that no verifier accepts --- an id not formed from its preimage, a parent
reference to nothing, an honest block's id copied onto other contents ---
lies outside the assumption altogether, so it cannot falsify it. That is
why the assumption is scoped to admissible blocks: claimed for every
signed block, \code{chain\_order} and collision resistance would be false
the moment an exposed key signed such junk, and every theorem below would
hold vacuously. Honest clocks
may run fast by up to $\sigma$; slow clocks cost only liveness.

One forward note for key rotation: with versioned keys, a leaked
\emph{retired} key is exposure of its seat --- as stated, it counts
against every window it can still verify, which without a pin is every
later one. Section~\ref{sec:results} recovers safety by validator
rules, with stolen keys of any version signing anything forever: mode 2
restricts this assumption to the blocks its validator admits, so a key
stops exposing a stamp once the schedule retires it; mode 3 rescopes it
to a per-key \emph{surface} defined there.

\begin{assumption}[Hash collision resistance, over blocks that
occur]\label{ass:hash}
For the headline theorem, collision resistance holds over the admissible
blocks that occur: any two admissible blocks (or the genesis) occurring in
the execution with equal ids are equal (\code{SigningExecutionOn.id\_inj}
over \code{SignedEver}). This is
classical collision resistance in the form a proof can consume, not a
weakening of it: plain injectivity of $H$ is false outright --- ids are
shorter than blocks, so colliding values exist in the type --- while
collision resistance says only that no adversary ever \emph{exhibits}
one. Scoping the domain to blocks that can actually occur is what turns
the second statement into the first; every block of a validated
certified chain is in that domain.
We call this \emph{hash injectivity}. Mode 2's theft theorem uses it in
exactly this form; for mode 3 and the untimed mode-2 form the domain is
blocks verifying under their \emph{declared} version
(\code{SignedDeclared}) ---
Section~\ref{sec:protocol} explains why \code{keyIndex} is in the
preimage. \quad\textit{(Lean: \code{SigningExecutionOn.id\_inj},
\code{SignedHashInjective}.)}
\end{assumption}

\begin{assumption}[Certificate grounding]\label{ass:cert}
A certificate that verifies carries a claim built from the deployment
genesis by folding in one block at a time, where each fold checks the link
rules, the density of newly matured windows, and the folded block's
admission: a verifying producer signature and a formed id. Deliberately minimal: it does \emph{not}
posit a valid chain behind the certificate --- a proved lemma
reconstructs the prefix chain from the fold steps, so the chain's
existence is a consequence, not a hypothesis.
\quad\textit{(Lean: \code{GroundedCert}.)}
\end{assumption}

\begin{assumption}[Clock]\label{ass:clock}
The verifier reads local time $\mathit{now}$, and real time is at most
$\delta$ slots ahead of it ($R \le \mathit{now} + \delta$). The verifier
accepts a tip when $\mathit{now} \le \mathit{tip.slot} + b$. Theorem~1
then applies at freshness $\varphi = b + \delta$, because
$R \le \mathit{now} + \delta \le \mathit{tip.slot} + b + \delta$.
Here $\delta$ covers clock skew and propagation combined. No relation
between $b$ and $n$ is needed: a larger $b$ tolerates slower delivery at
the price of a longer stretch of real time in which exposure counts
against each window (Assumption~\ref{ass:budget}).
\end{assumption}

\begin{assumption}[Honest delivery --- liveness only]\label{ass:delivery}
For the liveness theorem only: the honest producers of each newly maturing
window were live, and their blocks reached the current producer in time to
be built on. \quad\textit{(Lean: \code{HonestBlocksCover}.)}
\end{assumption}

\section{What the theorems guarantee}
\label{sec:results}

The results come in a fixed order. First, light-client safety
(Theorem~\ref{thm:lc}), proved in the timed model
(\code{timed\_light\_client\_safety}). Next, no valid chain runs ahead of
real time (Theorem~\ref{thm:forge}), the fact the proof of
Theorem~\ref{thm:lc} rests on. Then the rotation results per mode
(Section~\ref{sec:rotation}), and finally liveness.

Two adversaries appear, and it matters which one a result faces. The
\emph{exposure} adversary of Section~\ref{sec:assumptions} obtains keys
at arbitrary real times, by theft or by transient control, within the
budget. Theorems~\ref{thm:lc} and~\ref{thm:forge} face it with a stamp
counted as exposed whichever of its seat's keys is; mode 2's theft
theorem faces the same adversary with exposure read at the schedule, so
a stolen key exposes only the stamps it can still sign admissibly ---
which is what rotation buys. The \emph{key-stealing} adversary of
Section~\ref{sec:rotation} is charged through a per-slot census of
rented slots and stolen live versions instead; mode 3 and the untimed
mode-2 form face it.

Every result below is a Lean theorem, proved once against the model.
The light-client theorem (Theorem~\ref{thm:lc}) is carried to each
implementation by that implementation's soundness bridge
(Section~\ref{sec:impl}): \code{Rust.rust\_timed\_certified\_agreement}
for the Rust validator and \code{ts\_timed\_certified\_agreement} for the
TypeScript one. The Rust bridge also covers the validator's
key-discipline conjunct. Theorem~\ref{thm:forge}, the rotation-mode
results of Section~\ref{sec:rotation} and the liveness theorem are,
today, proved at model level. A guarded \code{\#print axioms}
per headline theorem is checked in (\code{Molt/Axioms.lean}), so the build
fails --- rather than a reviewer noticing --- if any proof ever assumes
more than classical logic.

\subsection{Light-client safety}
\label{sec:lc}

\begin{theorem}[Recent chains share their deep blocks; Lean
\code{timed\_light\_client\_safety}]\label{thm:lc}
Fix $n \ge 1$, a deployment genesis $G$, a freshness $\varphi$, an honest
clock lead $\sigma$, and a lookback $\ell$ with $\ell \ge n - 1$ and
$\ell \ge n + \fmax + \sigma + 1 - \quorum$. Suppose:
\begin{itemize}
\item a signing execution over the verifier's admissible blocks
  (\code{SigningExecutionOn}~$\mathit{Signed}$, where $\mathit{Signed}$ is the
  per-block check: the signature verifies and the id is formed):
  signing happens at real slots; while the key that verifies a stamp is
  unexposed, at most one admissible block with that stamp is ever signed
  under it (\code{honest\_once}); an admissible block is signed only once
  its parent is available (\code{chain\_order}); and collision resistance
  holds over admissible occurring blocks (\code{id\_inj}).
  Exposed keys are unconstrained: the first field speaks only of
  unexposed keys, the other two are cryptographic facts, and blocks the
  verifier rejects are not mentioned at all;
\item honest clocks lead real time by at most $\sigma$ slots, for
  admissible blocks (\code{HonestClockOn});
\item the exposure budget (\code{ExposureBounded}~$n$~$\ell$~$\varphi$): of
  any $n$ consecutive stamps $[v, v + n)$, at most $\fmax$ have their key
  exposed at some real slot in $[v - \ell,\ v + n + \varphi)$;
\item the verifier is at real slot $R$, and every admissible block it accepts
  was produced at some real slot no later than $R$ (\code{hbridge}: EUF-CMA
  plus causality);
\item there are two presentations, each consisting of a grounded certificate
  claim (\code{GroundedCert}) and a suffix that passes the validator:
  linkage to the claim (\code{hLink}), internal linkage (\code{hLinks}),
  density of maturing windows (\code{hDense}), and admissible blocks:
  verifying signatures and formed ids (\code{hSigned});
\item both tips are fresh in real time: $R \le \mathit{tip.slot} + \varphi$.
\end{itemize}
Then, taking any histories the two groundings attest (\code{GroundedHistory}),
the two full chains (attested history concatenated with suffix) agree at
every height $h$ with $h + n$ below both full lengths. Equivalently, they
agree on everything at least $n$ blocks below the lower tip.
\end{theorem}

Two remarks. First, the conclusion also covers a chain the node
\emph{produced} rather than received, because production ships nothing
that does not pass the same validator (Section~\ref{sec:rules}). Second,
the proof is the classic count with its timing made explicit. Take the
last full window of $n$ stamps below the lower tip: each chain holds at
least $\quorum$ blocks there, so at least $2\quorum - n \ge \fmax + 1$
stamps carry a block on both. Were every such pair distinct,
\code{honest\_once} would make each of those stamps' keys exposed at one
of the two signing times. Each signature was made by $R$, at most
$\varphi$ slots after the window ends, and --- by Theorem~\ref{thm:forge}
--- no earlier than $\ell$ slots before it starts. That is $\fmax + 1$
exposed stamps inside one budget stretch, which the budget forbids. So the
two chains share a block in that window, and the parent-id chain (id
formation plus collision resistance) makes them agree on everything below
it.

A third remark, on the lookback. Counting every exposure up to
$\varphi$ slots after a window --- however long before it --- removes the
clock hypothesis: \code{exposure\_agreement\_ever\_on} proves the full-chain
form of Theorem~\ref{thm:lc} under that cumulative budget
(\code{ExposureBoundedEver}) without \code{HonestClockOn}. The lookback is
what lets a signer that was compromised and then cleaned up stop counting
against later windows; the honest clock lead is its price. Under the
cumulative budget the no-early-signing bound also holds, as
$B.\mathit{slot} \le r + n + \fmax + \sigma + 1 - \quorum$ for
$B.\mathit{slot} \ge n - 1$ (\code{exposure\_no\_early\_signing\_ever\_on}).

\subsection{No chain runs ahead of real time}
\label{sec:timed}

The recency rule needs one more fact: an adversary must not be able to
build, early, a chain whose tip will look recent later. The id-formation
contract of Section~\ref{sec:protocol} appears as the field
\code{chain\_order}: a block can be signed only once its parent is
\emph{available}, because the parent's id preimage contains the parent's
signature, and predicting another party's signature on a known message
is exactly an EUF-CMA forgery. (This holds even for a deterministic
signature scheme: the value is determined, but computing it without the
key is the forgery the game forbids.)

\begin{theorem}[No valid chain runs ahead of real time; Lean
\code{no\_early\_signing}]\label{thm:forge}
Under the signing execution, honest clocks and exposure budget of
Theorem~\ref{thm:lc} (lookback condition included), let $c$ be a valid
chain from genesis $G$ whose non-genesis blocks are admissible and all
exist by some real slot, and let
$B \neq G$ be one of its blocks. If $B$ was signed at real slot $r$, then
\[
B.\mathit{slot} \;\le\; r + \ell .
\]
With the least lookback the condition allows --- for $n = 3\fmax + 1 \ge 4$
and $\sigma \le \fmax - 1$ that is $\ell = n - 1$ --- no valid chain gets
more than $n - 1$ slots ahead of real time.
\end{theorem}

The proof is one count, by induction along the chain. Suppose $B$ was
signed more than $\ell$ slots before its stamp. Its signer was then not
honest, whose clock leads by at most $\sigma \le \ell$, so $B$'s key was
exposed at $r$. The window of $n$ stamps ending at $B$ is matured at $B$,
so it holds at least $\quorum$ blocks, all of them $B$ or its ancestors,
all signed no later than $r$ (\code{chain\_order}). An honest one is
stamped at most $r + \sigma$. An exposed ancestor is, by induction, not
signed early either, so its exposure falls inside the window's budget
stretch and costs the budget one stamp. If $B$'s own exposure falls
inside that stretch too, the honest stamps cover too little of the window
to reach $\quorum$. If it does not, $B$ was signed at least $n$ slots
before the window starts, so the window holds no honest block and not the
genesis block, and its budget leaves too few blocks for $\quorum$. (For
$n = 1$ the same count runs on the window of $B$'s predecessor stamp.)

What this rules out, and what it does not. A patient adversary can still
bank signatures into a long fork, by exposing keys long after the stamps
it signs. Such a fork is stale: late exposure does not count against old
windows, but it cannot produce recent blocks either, because recent
windows have their own budget. A stale fork is exactly what the recency
check refuses. Theorem~\ref{thm:forge} rules out the converse: a fork
prepared early so as to look recent later.

\subsection{Key rotation}
\label{sec:rotation}

The protocol's key-rotation machinery was defined with the protocol
itself: the in-band version (\code{keyIndex}, Section~\ref{sec:protocol})
and the monotone rule (\code{keyMonoOk}) are already part of the one
chain validity. What this subsection adds is the last conjunct a
deployment fixes --- the \emph{pin}: a bound on the version a block may
declare (a floor in modes 1--2, exact per grid window in mode 3), which is
what makes a retired key \emph{dead} rather than merely old. The three
modes differ in exactly one thing ---
\textbf{what computes the pin}: the chain itself (mode 1), slot
arithmetic (mode 2), or a roster-wide counter read off the signatures
(mode 3). That one choice determines everything a light client must hold,
so we state each mode as: the rule, what the operator must do, what the
client must do, and what theorem you get.

\paragraph{The adversary.} Modes 2 and 3 face the
key-stealing adversary: a stolen key of \emph{any} version signs
anything, at any slot, forever --- no forward
security~\cite{forwardsec} is assumed. Mode 2 meets it with the
assumptions of Theorem~\ref{thm:lc}, read at the schedule (below).
Mode 3, whose pin depends on the chain, and the untimed mode-2 form meet
it through a per-mode \emph{signature surface}: the bundle of custody and
unforgeability assumptions the theorems consume, the analogue of
Assumption~\ref{ass:sig}, under which the honest-slot uniqueness
guarantee holds only for non-stolen, non-rented entries. Their budget
(refining Assumption~\ref{ass:budget}) is two rates per $n$-slot window.
The first: $\rho$ \emph{rented} slots --- slots the adversary fully
controls for that one slot and then loses, the shortest transient
exposure of Assumption~\ref{ass:budget}. A rented slot yields signatures
stamped for that slot only: a rented device made to sign blocks stamped
for other slots has had that key version stolen, and is counted in $T$.
This is the reading the per-mode signature surfaces encode. The second:
$T$ producers with a \emph{then-live} key stolen --- a stolen version
at-or-above the floor where it is counted. Together:
$\rho + T \le \fmax$. The key-stealing corruption predicate
(\code{badSched} in the untimed mode-2 form; mode 3 reads it at the
roster's counter) counts a slot bad if it is rented \emph{or} some
version at-or-above the one in force for its producer is stolen ---
which builds \emph{healing} into the definition: a stolen key stops
counting the moment its retirement takes force. Each consumes its own
rescoping of the surface --- the \emph{scheduled surface} (untimed
mode 2), the \emph{core surface} over the pin-free validator (mode 3)
--- and each rescoping includes \emph{version independence}: stealing
version $j$ yields nothing about $j' > j$, which key-insulated
signatures~\cite{keyinsulated,keyinsulatedsig} provide off the shelf. Hash injectivity
(Assumption~\ref{ass:hash}) is consumed alongside, over declared-version
signed blocks (\code{SignedDeclared}).

\begin{assumption}[Honest custody --- per-key surfaces]\label{ass:rotation-honest}
For mode 3 and the untimed mode-2 form: fix the rotation mode, a stamp $s$ whose seat is not rented, and a key version $j$ of its producer $i$ that has never been stolen. Every block stamped $s$ that verifies under version $j$ was signed by $i$'s own signing device, and that device signs at most one block per stamp, ever (a persisted signing record), declaring the version the mode's rule assigns to that stamp. The device's clock may lead real time by up to $\sigma$ slots, as in Assumption~\ref{ass:sig}; nothing here asks for an exact clock. This constrains the honest, uncorrupted signing path only: the instant a delegate key is rented or stolen, the adversary of Section~\ref{sec:rotation} is free of it, exactly as elsewhere in this paper (a stolen key signs anything, at any slot, forever). It is stated directly, as an operational assumption, rather than derived from the signing execution of Theorem~\ref{thm:lc}: \code{SchedUnforgeable}/\code{SchedCoreUnforgeable} are stated per individual key version, and one real slot can host several eligible key versions of which only one verifies --- so one key's safety does not certify the slot's (machine-checked: \code{badSched\_single\_key\_safe\_not\_enough}).
\end{assumption}

\emph{Without theft, formally.} A \emph{lost} key (destroyed, not
exfiltrated) gives the adversary nothing; the loser rotates and the
validator already supports it. With nothing stolen the corruption the
theorems count collapses to rent alone in every mode
(\code{badKeyrot\_lossOnly}, \code{badSched\_lossOnly}; mode 3 inherits
mode 2's count), and Theorem~\ref{thm:lc} applies to each mode's
validator unchanged: \code{keyrot\_loss\_agreement} (mode 1,
Theorem~\ref{thm:keyloss} below), \code{sched\_loss\_agreement}
(mode 2), \code{lockstep\_loss\_agreement} (mode 3). No anchor, no
cadence, any pin: a genesis-and-clock client gets the core guarantees of
Sections~\ref{sec:lc}--\ref{sec:timed}.

\paragraph{Mode 1: reactive rotation.} \emph{The rule.} A producer
rotates whenever it chooses, by signing under a higher version. The pin
is read off the chain being validated: seat $i$'s version \emph{in force}
at slot $s$ (\code{inForce}) is the highest version $i$ has declared in
the chain's \emph{confirmed prefix} --- blocks at least $n$ slots old ---
and no block may declare less (\code{validSignedChainK$'$}). Why $n$
slots: a rotation must take effect at one agreed moment on every chain,
``the announcing block is confirmed'' is that moment, and $n$ is exactly
the depth at which all accepted chains provably agree on the confirmed
prefix (the Lean proof takes this depth as a parameter and requires
only that it be at least $n$). At the prototype's parameters a rotation takes force in under
half a minute. 

\emph{What it guarantees: recovery from key loss.} A seat that loses its
delegate key --- destroyed, held by nobody --- signs its next block under
a higher registered version, and production resumes
(\code{liveness\_produce\_blockK}). The replacement version must
already be in the deployment's key registry, for instance delegated from
the seat's cold root; the verifier never learns a new key from the chain.
Loss changes none of the hypotheses of Theorem~\ref{thm:lc}, so mode 1
inherits it directly.

\begin{theorem}[Mode 1 under key loss; Lean
\code{keyrot\_loss\_agreement}]\label{thm:keyloss}
Fix $n \ge 1$, a deployment genesis $G$, and the signing execution,
honest clocks and exposure budget of Theorem~\ref{thm:lc}
(\code{SigningExecutionOn}, \code{HonestClockOn}, \code{ExposureBounded},
lookback condition included), where a block is admissible when it
verifies under its declared key version (\code{SignedDeclared}) and its
id is formed. The verifier is
at real slot $R$, and every admissible block was produced at some real
slot no later than $R$ (\code{hbridge}). Take two signed chains that both
pass the mode-1 validator (\code{validSignedChainK$'$}, at any
confirmation depth), with formed ids (\code{hFormed}), both start at $G$,
and both have tips recent in real time
($R \le \mathit{tip.slot} + \varphi$). Then they agree at every height $h$
with $h + n$ below both lengths.
\end{theorem}

\emph{What it does not guarantee: key theft.} Mode 1 is not safe against
key theft, and we claim nothing there. The pin keeps a retired key dead
only on chains that contain its retirement; a fork that branches before
the retirements is under no such constraint, and on it the old versions
are still current. Producers stop guarding retired keys, so an adversary
that later steals the retired keys of a quorum of seats can build a
dense fork from a point before their rotations, stamp its blocks so that
its tip is recent, and present it to a client holding the genesis and a
clock. The fork passes the mode-1 validator and the recency check, and it
disagrees with the real chain arbitrarily deep. No fault budget on the
real chain excludes it, since the stolen keys were retired there.
Excluding such a fork needs a client that holds a recent block of the
real chain and keeps re-verifying, with a theft budget charged to every
window a theft can reach; we do not state that rule as a result. A
deployment that must survive key theft runs mode 2 or mode 3.

\paragraph{Mode 2: scheduled rotation.} \emph{The rule.} The version in
force at slot $s$ --- modes 2 and 3 call it the \emph{generation} --- is
a protocol constant
$\mathit{gen}(s) = \lfloor s/E \rfloor$ for an era length $E$: a pure
function of the slot, identical on every chain, pinned at the verifier
and never peer-supplied. The validator adds one check:
$\mathit{gen}(\mathit{slot}) \le \mathit{keyIndex}$
(\code{validSignedChainSched}). Because the pin reads nothing off any
chain, a fork \emph{cannot disagree} with the real chain about which keys
are retired --- the fake-past construction that defeats mode 1 is
structurally impossible.

\emph{What it guarantees: theft safety, from Theorem~\ref{thm:lc}
itself.} Mode 2 needs no assumption beyond those of
Theorem~\ref{thm:lc}; it only reads them at the schedule. Write
$\mathit{Controlled}(i, r)$ for ``seat $i$ is controlled by the adversary
at real slot $r$'' and $\mathit{Stolen}(i, j, r)$ for ``someone other
than seat $i$'s honest holder can sign under its version-$j$ key at real
slot $r$''; both are arbitrary in time. Stamp $s$, with producer
$i = s \bmod n$, is \emph{exposed} at $r$ when $\mathit{Controlled}(i, r)$
holds or $\mathit{Stolen}(i, j, r)$ holds for some $j \ge \mathit{gen}(s)$
(\code{schedExposed}). A block is \emph{admissible} when it verifies under
its declared version, that version is at least $\mathit{gen}$ of its
slot (\code{SchedAdmissible}), and its id is formed --- the per-block
content of the mode-2 validator and its hashing layer. Assumption~\ref{ass:sig} is needed only for admissible blocks
(\code{SigningExecutionOn}, \code{HonestClockOn}): a block signed under a
key already retired at its stamp falls outside it, since no validator
accepts that block.

\begin{theorem}[Mode 2 under key theft; Lean
\code{sched\_exposure\_agreement},\hfil\penalty0
\code{sched\_exposure\_certified\_agreement}]\label{thm:sched}
Fix $n \ge 1$, a deployment genesis $G$, a schedule $\mathit{gen}$, and any
control and theft histories $\mathit{Controlled}$ and $\mathit{Stolen}$.
Assume the signing execution and honest clocks of Theorem~\ref{thm:lc}
for admissible blocks, over the mode-2 exposure; the exposure budget of
Theorem~\ref{thm:lc} (\code{ExposureBounded}~$n$~$\ell$~$\varphi$) over
the same exposure, lookback condition included; and that the verifier is
at real slot $R$ and every admissible block was produced at some real
slot no later than $R$ (\code{hbridge}). Take two signed chains that both
pass the mode-2 validator (\code{validSignedChainSched}), with formed ids
(\code{hFormed}), both start at $G$, and both have tips recent in real time
($R \le \mathit{tip.slot} + \varphi$). Then they agree at every height $h$
with $h + n$ below both lengths. The same holds for two certificate
presentations in the form of Theorem~\ref{thm:lc}, with claims grounded
under the scheduled grounding (\code{GroundedCertSched}: each fold adds
the pin check) and admissible suffix blocks: any histories the two
groundings attest (\code{GroundedHistorySched}), each concatenated with its
suffix, agree at every height $h$ with $h + n$ below both full lengths.
\end{theorem}

Both forms are instances of one strengthening of the core, proved once:
Theorems~\ref{thm:lc} and~\ref{thm:forge} hold with
Assumption~\ref{ass:sig} restricted to any per-block predicate that every
non-genesis block of both chains satisfies
(\code{exposure\_agreement\_on}, \code{exposure\_no\_early\_signing\_on},
\code{exposure\_certified\_agreement\_on}); the unrestricted theorems are
implied by the case of the trivial predicate, but that case is not
an assumption a deployment can make: claimed for every signed block,
the execution fields fail as soon as an exposed key signs a block with a
dangling parent reference. The restriction is what the
rotation pin pays for: the exposure is per stamp, over every version that
can sign that stamp admissibly, so the gap between one key's safety and a
slot's safety that the per-key surfaces below must work around does not
arise.

\emph{What the budget charges.} The budget is
Assumption~\ref{ass:budget}'s, and the schedule makes it fine-grained. A
stolen version $j$ exposes only stamps whose generation is at most $j$.
Since theft persists, it is charged against every window whose stretch of
real time reaches the theft, until the schedule passes $j$: from about
the theft's real time to the end of era $j$, and never against a window
whose stretch ended before the theft. Healing is a consequence of the
exposure, not a definition. No erasure is required: a retired key, stolen
or not, exposes nothing at any stamp it can no longer sign. A key stolen
\emph{early} --- derived ahead of its era and taken --- is a stolen
version at or above the floor of every stamp until its era ends, and is
charged accordingly; a stolen \emph{cold root} is a theft of every later
version of its seat at once, so that seat counts as exposed in every later
window, like a Byzantine seat. Key-provisioning discipline is therefore
not an extra assumption: it is how an operator stays within the budget the
theorem consumes. One modelling rule remains: $\mathit{Stolen}$ must list
every version the adversary can sign under. \emph{Version independence}
--- stealing version $j$ yields nothing about $j' > j$, which
key-insulated signatures~\cite{keyinsulated,keyinsulatedsig} provide off the shelf --- is
what keeps that list as short as what was actually taken.

\emph{Operator duties.} Keep one offline cold root per seat and delegate
each generation's signing key just in time for its era; then a theft costs
only the rest of the live era. The era length $E$ is a deployment dial:
longer eras mean rarer cold-root ceremonies but a longer budget charge per
stolen live key. A retired key is harmless \emph{by rule}, not by secrecy;
separately, a chain whose tip declares a generation whose era ended more
than $n$ slots ago fails the recency check
(\code{sched\_oldkey\_fork\_stale}). \emph{What breaks first:} (i) a
peer-supplied schedule --- the pin must be a verifier constant, or a zero
schedule from the network readmits every retired key; (ii) theft of the
\emph{current} generation, charged until its era ends --- size the budget
and $E$ together; (iii) a stolen cold root, which costs its seat for
good.

\emph{Without a genesis.} An earlier, untimed form keeps the full-chain
client genesis-free, deriving genesis agreement instead of assuming it
(\code{sched\_recent\_tip\_ancestor\_agreement\_horizon},
\code{sched\_recent\_tip\_ancestor\_mem\_horizon},
\code{sched\_recent\_genesis\_agreement\_horizon}): at equal tip heights
two accepted chains agree on the block $n$ below each tip, and at unequal
heights the lower chain's $n$-deep ancestor is a block of the taller
chain, at least $n$ deep there; and two accepted chains each longer
than $n$ blocks share their genesis. It consumes the scheduled per-key
signature surface (\code{SchedUnforgeable},
Assumption~\ref{ass:rotation-honest}) and the untimed census budget of
the adversary paragraph above, from a horizon $H$ on
(\code{ByzantineBoundedFrom}~$H$~$n$~(\code{badSched}~\dots), with side
conditions $H + n \le \mathit{tip.slot} + 1$), and its census counts a
stolen not-yet-retired version from the start of time, so early theft and
cold-root custody must be assumed rather than charged. Its certificate
form is \code{sched\_recent\_certified\_suffix\_agreement\_horizon}. We
keep it in the artifact for a client that cannot hold a genesis;
Theorem~\ref{thm:sched} supersedes it otherwise.

\paragraph{A realistic mode-2 deployment, end to end.} At the ceremony:
each seat generates its cold root offline and the deployment publishes
$\mathit{gen}(s) = \lfloor s/E \rfloor$ with, say, $E$ = one week of
slots. Weekly, each seat's cold root signs one delegation for the
coming generation and goes back in the safe; the live signing key never
predates its era by more than the provisioning margin. Clients are
wallets holding the genesis, a clock and $\mathit{gen}$; they sleep for
months and sync on wake. What must then be true of the world:
Assumption~\ref{ass:sig} for admissible blocks, hash collision resistance
(Assumption~\ref{ass:hash}), the standing clock assumption
(Assumption~\ref{ass:clock}), and the windowed budget of
Assumption~\ref{ass:budget} over the mode-2 exposure. A stolen live key
is charged for at most the rest of its week plus the window's stretch,
and that is what the budget is sized against. Through
Theorem~\ref{thm:sched} and the parent-id chain
(\code{same\_block\_same\_prefix}), any two accepted recent chains agree
on everything at least $n$ blocks below the lower tip, whatever happened
while the wallet slept. A certificate client gets the same guarantee
(\code{sched\_exposure\_certified\_agreement}).

\paragraph{Mode 3: free-cadence lockstep.} \emph{The rule.} No schedule
anywhere. The roster advances a single shared generation counter at
moments of its own choosing --- an incident, a governance decision ---
and validity enforces only the \emph{discipline} of the switch
(\code{noMixing}): within each grid window (slots $[Wn,\,Wn{+}n)$; the
fixed grid, not a sliding window) all blocks declare one generation, and
the counter never decreases across grid windows
(\code{validSignedChainLock}). The verifier never learns the cadence; it
reads each grid window's generation off the signatures themselves.

\begin{theorem}[Lockstep safety; Lean
\code{lockstep\_client\_safety}]\label{thm:lock}
Under the lockstep validator and the lockstep hypothesis package (the
roster tracks one monotone counter; honest signatures declare their
grid window's value; the genesis declares the counter's initial value; hash
injectivity; the same $\rho$/$T$ budget, assessed at the
one-grid-window-lagged counter over the pin-free core validator's signature
surface): the anchor-free agreement of the untimed mode-2 form ---
at equal tip heights, on the block $n$ below each tip; at unequal
heights (\code{lockstep\_recent\_tip\_ancestor\_mem}), the lower chain's
$n$-deep ancestor is a block of the taller chain, at least $n$ deep there. The client holds even less than
mode 2's --- genesis and a clock, with no schedule to store at all.
\end{theorem}

The route is a machine-checked \emph{pinning} theorem
(\code{lockstep\_declares\_rosterGen}): on any accepted recent chain,
every matured grid window declares exactly the roster's counter --- a
grid window at any other generation would need a quorum of adversarial
blocks, which
a $\rho + T \le \fmax$ adversary cannot supply. Pinning makes every
accepted lockstep chain pass mode 2's validator at the (one-grid-window-lagged)
counter, so the untimed mode-2 theorems transport.

\begin{theorem}[Lockstep safety, per-generation census; Lean \code{lockstep\_client\_safety\_gen},\hfil\penalty0
\code{lockstepGen\_recent\_tip\_ancestor\_mem}]\label{thm:lock-gen}
Under the lockstep validator (\code{validSignedChainLock}) and the
per-generation package --- honest signatures declare their grid window's
counter; hash injectivity; the pin-free core signature surface; rent at most
$\rho$ per window and, for every generation $j$ separately, at most $T$ seats
whose generation-$j$ key is stolen; $\rho + T \le \fmax$; recency
$\mathit{now} \le \mathit{tip.slot} + \Delta$ for both tips; and chain length
$2n < \mathit{length}$ (\code{hLong}, for both chains at equal heights;
for the lower chain at unequal heights) --- any two accepted chains agree:
at equal tip heights (\code{hTipHeight}), on the block $2n$ below each tip
(\code{lockstep\_client\_safety\_gen}); at unequal tip heights, the lower
chain's $2n$-deep block is a block of the taller chain, at least $2n$ deep
there (\code{lockstepGen\_recent\_tip\_ancestor\_mem}).
Through the parent-id chain, two such chains both longer than $2n$ blocks
agree on their whole common prefix down to height $0$: no shared genesis
is assumed
(\code{lockstepGen\_recent\_genesis\_agreement}).
\end{theorem}

The route is a second pinning theorem that needs neither induction nor genesis (\code{lockstep\_window\_declares\_rosterGen}): the last grid window matured below the lower tip declares one generation on both chains, so the counting argument runs on a window whose corruption reading is at a single generation (\code{window\_shared\_prefix}). The unmatured tip window cannot be used, and that is what costs depth $2n$ --- exactly $n + ((t{+}1) \bmod n)$ positions for lower tip slot $t$ (\code{lockstepGen\_shared\_prefix\_sharp}). The two packages are incomparable in general; the cumulative one delivers the per-generation one when the counter skips no generation (\code{LockstepPackage.toGen}).

\emph{Operator duties.} Two, and they are the mode's real price:
\emph{lockstep coordination} --- all honest signers track the one counter
(a split here is an operational failure the theorems do not cover) ---
and \emph{real erasure} of retired key material at each switch. Its budget credit is now formal: with theft time-stamped, erasure is the hypothesis that no generation is stolen after the roster has moved past it (\code{ErasureTimedLock}), and it is \emph{consumed} --- it turns the census a deployment attests while a generation is live into the per-generation census Theorem~\ref{thm:lock-gen} needs (\code{genBound\_of\_preRetirementBound}; end to end, \code{lockstep\_client\_safety\_timed}). A deployment budgets $T$ seats lost \emph{per generation, while it is live}, not $T$ across all generations. In mode 2 neither erasure nor provisioning discipline is a hypothesis: the schedule makes a retired key harmless, and the timed budget charges early theft directly (Theorem~\ref{thm:sched}); mode 3, with no schedule, needs erasure instead. Erasure
that is not real is silent to the protocol: keys held in an HSM or only
ever in RAM are defensible; flash storage, which does not erase by
overwriting, is not.

\paragraph{A realistic mode-3 deployment, end to end.} Each seat signs
from a hardware module whose retired generations are destroyed on
switch; the genesis block declares the counter's initial value; switches
are agreed out of band --- an incident response, a governance action ---
and every honest signer moves at the agreed grid-window boundary. Clients
hold genesis and a clock, nothing else, and sleep arbitrarily long; on
wake they sync from a certificate plus an $O(n)$ suffix, as in modes
1--2. The certificate carries the claim and the tip generation; for
$n \ge 2$ the generation is determined by the claim's own tail
(\code{groundedCertLock\_gen\_of\_tail}, \code{groundedCertLock\_gen\_unique}),
so the certificate attests nothing beyond the plain claim, and only
$n = 1$ needs the extra counter --- which must then be attested with
the claim: an unauthenticated counter readmits a stolen retired-generation
key. What must then be
true of the world: the core surface over the pin-free validator and hash
injectivity (the adversary paragraph's rescopings); the standing clock
assumption (Assumption~\ref{ass:clock}); the coordination discipline above (a
silent split among honest signers is an operational failure outside the
theorems); and \emph{either} the same $\rho$/$T$ budget as the untimed
mode-2 form, read globally, assessed at the one-grid-window-lagged counter
(Theorem~\ref{thm:lock}, depth $n$), \emph{or}, at confirmation depth
$2n$, rent at most $\rho$ per window and, for every generation $j$, at
most $T$ seats whose generation-$j$ key is stolen while $j$ is live
(Theorem~\ref{thm:lock-gen}; end to end with time-stamped theft,
\code{lockstep\_client\_safety\_timed}).
Conclusion, machine-checked (\code{lockstep\_client\_safety} and
\code{lockstep\_recent\_tip\_ancestor\_mem}, via the pinning theorem
\code{lockstep\_declares\_rosterGen}; and at the certificate presentation under
the cumulative budget at depth $n$,
\code{lockstep\_recent\_certified\_suffix\_agreement}, and under the
per-generation census at depth $2n$,
\code{lockstepGen\_recent\_certified\_suffix\_agreement}, with the
certificate's own counter pinned to the roster's once its grid window
matures, \code{lockstep\_cert\_gen\_pinned}): any two accepted chains agree ---
at equal heights on the block $n$ below each tip (or $2n$ under
the per-generation census), at unequal heights the
lower chain's confirmed block lying on the taller --- hence, through the
parent-id chain (\code{same\_block\_same\_prefix}), on the lower chain's
whole confirmed prefix --- with no schedule anywhere and no anchor.

The modes at a glance and their headline theorems are collected in
Section~\ref{sec:opmodel}.

\paragraph{Honest scope.} Three bounds on what this subsection claims.
First, mode 1 carries no theft guarantee (see above).
Second, the mode-2 and mode-3 results are model-level: the imported Rust
validator carries the mode-1 floor check (bridged in \texttt{Rust/BridgeK.lean})
but neither the schedule pin nor the lockstep rule.
Third, mode 3's signature surface, and that of the untimed mode-2 form,
are not derived from the timed model: as witnessed by
\code{badSched\_single\_key\_safe\_not\_enough} and
\code{badKeyrot\_single\_key\_safe\_not\_enough}, single-key safety does not
imply slot safety, so their honest custody is stated directly as an
operational assumption (Assumption~\ref{ass:rotation-honest}). Mode 2's
theft theorem has no such surface: it consumes Assumption~\ref{ass:sig}
itself, restricted to admissible blocks.

\subsection{Liveness, and the slot time}
\label{sec:liveness}

Safety bounds what a verifier can be made to accept; liveness is about
progress, and its proofs use only the structural core of the paper's
single chain validity (Section~\ref{sec:protocol}) --- no rotation
hypothesis enters.

\begin{theorem}[Liveness; Lean \code{production\_liveness},
\code{global\_liveness}]\label{thm:live}
\emph{Local:} with at most $\fmax$ slots per $n$-slot window whose
producers fail (\code{ByzantineBounded}) and honest delivery
(Assumption~\ref{ass:delivery}), when a slot
arrives that belongs to an honest producer holding an accepted chain,
production succeeds --- the extension passes the same validator everyone
runs. \emph{Global:} a synchronous honest run within the budget (over
the run's span; each honest block reaches the next honest producer
within its slot) is
validator-accepted with one block per honest slot, so the chain reaches
any height and every block eventually becomes $n$-deep.
\end{theorem}

The fork-choice rule --- adopt the valid chain with the highest tip ---
is what makes honest blocks coalesce into the single chain the global
form assumes; a fully multi-chain network model is future work.

One quantity remains: the slot duration $\tau$. It is set by prover
economics, not network latency: the slot must be long enough for the
proving pipeline to keep its uncovered suffix under $n$ blocks forever. The Lean
development proves a threshold that is sufficient, and necessary over
the pipeline's steady states
(\code{slot\_time\_sufficient}, \code{slot\_time\_necessary}); we do not
reproduce the formula here --- Section~\ref{sec:impl} instantiates it at
the measured proving costs.

\section{Implementation and measurements}
\label{sec:impl}

\paragraph{One source, three instantiations.} The reference
implementation is one Rust file, \code{rust/src/lib.rs}, written in the
fragment that the Charon + Aeneas toolchain translates to Lean~4. From
it:
\begin{itemize}
\item \emph{The native validator} is imported into Lean and proved sound
  against the model: if the extracted code accepts a chain, the semantic
  predicate holds (\code{rust\_valid\_chain\_sound} and companions).
  Soundness is the only direction safety needs --- completeness is a
  liveness concern --- and it is what carries Theorem~\ref{thm:lc} to the shipped code
  (\code{Rust.rust\_timed\_certified\_agreement}). Id formation is
  checked by the hashing layer, not by the extracted validator, so that
  theorem and its TypeScript twin take it as a hypothesis on the presented
  blocks (\code{hFormedS}) and as part of each grounding fold. The bridges cover the
  validator's structural core and its key-discipline conjunct
  (\code{rust\_valid\_chain\_sound}, \code{rust\_valid\_chain\_k\_sound});
  the three rotation \emph{pins} of Section~\ref{sec:rotation} are,
  today, stated and checked at model level.
\item \emph{The circuit} instantiates the same source's backend-generic
  validator (\code{valid\_chain\_be}) over plonky2, so the recursive
  certificate proves exactly the predicate that was verified
  (\code{Rust.valid\_chain\_be\_validChain}); the per-gadget faithfulness of
  the circuit backend is the one remaining trust assumption, stated
  rather than proved.
\item \emph{A TypeScript twin} (\code{moltPetit.ts}) reaches Lean through
  an independent compiler (Thales) with its own soundness bridge, which
  carries Theorem~\ref{thm:lc} to it (\code{ts\_timed\_certified\_agreement}); the
  vendored emission hand-corrects three recorded emitter bugs
  (\code{tools/thales-reemission/}), a hand-audited step in this
  pipeline's trusted base. It
  shows the methodology is not tied to one language, and keeping two
  implementations honest against one model is what kept the model
  abstract.
\end{itemize}
The trusted base is exactly: the Lean kernel; the pinned translators
(Charon + Aeneas, Thales together with the recorded hand-corrections to
its emission, the circuit gadgets); the hashing layer's id recomputation;
and the faithfulness of the named
assumptions to the deployed primitives. Crypto enters only as
theorem hypotheses, and every handle (keys, signatures, certificates,
payload commitments) is opaque to the protocol --- which is also why no
numeric-model axiom of either host language reaches any theorem.

\paragraph{Measurements.} The prototype bundles the consensus rule with
a minimal application --- a committed key set evolving by signed
rotations --- on plonky2~\cite{plonky2} over its native 64-bit field
(Goldilocks) with Poseidon hashing~\cite{poseidon}, signing with
SPHINCS+~\cite{sphincsplus,fips205} over Poseidon (parameter set
\code{128s}). All numbers are wall-clock on one Intel i5-13500 (14 cores,
20 threads). ``Circuit rows'' is the standard size measure of a statement
handed to the proof backend --- proving cost grows with rows --- and the
\emph{recursion floor} is the fixed minimum cost of verifying any
subproof inside another proof.

\begin{table}[t]
\centering
\small
\begin{tabular}{lrr}
\toprule
operation & time & circuit rows \\
\midrule
consensus validator alone (8-block chain) & $0.035$\,s & $2^{9}$ \\
one in-circuit signature verification (subproof) & $0.36$\,s & $2^{13}$ \\
verifying one subproof (recursion floor) & $0.19$\,s & $2^{12}$ \\
certificate, 8 signatures inline & $6.6$\,s & $2^{16}$ \\
certificate, 8 signatures as recursive subproofs & $3.0$\,s & $2^{15}$ \\
2048 subproofs, balanced tree (critical path) & $5.0$\,s & $2^{13}$/node \\
\bottomrule
\end{tabular}
\caption{What the prototype measures, and why it means the design is
viable: the verified consensus rule is ${<}2\%$ of any signature-bearing
certificate (row 1 vs.\ rows 4--6); post-quantum signatures dominate and
recurse well --- a signer's subproof is verified at the recursion floor,
and a balanced tree keeps latency logarithmic in approvals per slot
($5.0$\,s at $2048$, against ${\approx}770$\,s folding sequentially)
while hardware grows linearly.}
\label{tab:bench}
\end{table}

Two readings matter for deployment. First, \emph{consensus is free}: the
machine-checked validator adds no appreciable proving cost on top of
whatever the application already pays for signatures. Second, the
measured proving costs feed the slot-time rule of
Section~\ref{sec:liveness}, which at this hardware and $n = 7$ recommends
$\tau \approx 4$\,s --- the figure used throughout the paper. Where a
deployment tolerates stateful signing, XMSS~\cite{xmssrfc} over the same
Poseidon gadgets shrinks the in-circuit verifier from $2^{13}$ to $2^{10}$ rows
and speeds full certificates by roughly $10\times$, at the cost of
one-time keys whose bookkeeping the consensus layer already maintains.

\section{Related work}
\label{sec:related}

\paragraph{Density in proof-of-stake.} The window-density rule's nearest
relatives are in the Ouroboros family, where density is a first-class
quantity: Genesis~\cite{genesis} resolves deep forks by comparing growth
just after the last common block (a bootstrap that replays history), and
Samasika~\cite{samasika} --- Mina's consensus --- carries a constant-size
minimum-window-density summary, the closest prior pairing of density with
succinctness. The structural difference: in every Ouroboros variant
density drives \emph{fork choice}, so finality is probabilistic and
depth-dependent, while Molt Petit makes density a hard \emph{validity
predicate}, giving deterministic finality at a fixed depth of $n$ blocks.
Molt Petit also fixes an equally weighted roster with the public
round-robin leader $s \bmod n$, where the Ouroboros line is
stake-weighted, with private VRF-based election from
Praos~\cite{praos,ouroboros} onward --- and its security argument is a
machine-checked artifact, not a hand-written analysis.

\paragraph{Succinct certificates and light clients.} The constant-size
recursive certificate is anticipated by Mina~\cite{mina}, whose
recursively composed SNARK attests the \emph{whole} history; Molt Petit's
certificate is instead actionable only under the recency rule ---
recency-scoped rather than whole-history agreement.
Plumo~\cite{plumo} syncs ultralight clients by SNARK-proved committee
hand-offs, with a pen-and-paper argument over a committee-weighted set.
The proof-of-work line is probabilistic by nature:
NiPoPoWs~\cite{nipopow} via rare superblocks, FlyClient~\cite{flyclient}
via random sampling; both certify accumulated \emph{weight}, where Molt
Petit's certificate is a single recursive proof of a deterministic
validity predicate --- constant-size and sampling-free.

\paragraph{Deterministic-finality BFT.} Molt Petit shares with
Tendermint~\cite{tendermint} and HotStuff~\cite{hotstuff} a fixed
participant set, rotating leaders, and deterministic finality --- by the
opposite mechanism. They finalize through explicit quorum-certificate
\emph{voting}; Molt Petit exchanges no votes at all: the same $2n/3$
threshold is enforced by counting blocks in a window, and the light
client's constant-size object is a recursive proof of that counting rule,
not an aggregate of votes.

\paragraph{Long-range attacks, key evolution, weak subjectivity.} Keys
that leak after the fact are the classic long-range threat to
proof-of-stake and permissioned chains alike~\cite{stakebleeding}. Two
remedies dominate. Key-evolving (forward-secure)
signatures~\cite{forwardsec}, as in Ouroboros Praos~\cite{praos}, make an
old key useless once its holder erases it; weak
subjectivity~\cite{weaksubj} instead has a client that was offline too
long obtain a recent checkpoint from a trusted source. Molt Petit's
recency rule is the weak-subjectivity condition made a theorem
hypothesis, but on the tip's age as read from the client's own clock
rather than on a trusted checkpoint: a client may sleep arbitrarily long
and needs no trusted party on waking, and the windowed budget states
exactly which leaks a fork can still exploit --- only those that recency
rejects (Theorems~\ref{thm:lc}--\ref{thm:forge}). Against theft, mode 2's
public schedule does the work of key evolution without erasure --- a
retired key is dead by the validator's rule, not by its destruction ---
and mode 3 is the erasure-based variant with a free cadence.
Key-insulated signatures~\cite{keyinsulated,keyinsulatedsig} supply the
version independence both rotation modes rely on.

\paragraph{Machine-checked consensus.} Velisarios~\cite{velisarios}
(BFT in Coq), Verdi~\cite{verdi} (verified Raft, extracted OCaml), and
IronFleet~\cite{ironfleet} (Paxos-based replication in Dafny) all prove
running consensus code correct. Molt Petit differs in the object --- a
vote-free density consensus carrying a recursive light-client certificate
--- and in the pipeline's \emph{direction}: those systems refine or
extract downward from a prover-internal model, while Molt Petit
\emph{imports} the deployed sources into Lean~\cite{lean,mathlib} via
Charon + Aeneas~\cite{aeneas} and Thales~\cite{thales}, so the theorems
are about the already-runnable artifact. The trusted translation step
does not disappear --- the front-ends stand exactly where extraction
stands in theirs --- but it points from the shipped code into the prover
rather than out of it.

\section{Limitations and future work}
\label{sec:limitations}

\paragraph{What the design gives up.} Four properties of the surveyed
systems are deliberately traded away, and a deployment should know it is
buying their absence:
\begin{itemize}
\item \emph{Dynamic participation.} The roster is fixed and closed: no
  reconfiguration, no membership churn. Signing keys rotate
  (Section~\ref{sec:rotation}); seats do not.
\item \emph{Stake.} The set is equally weighted, the leader schedule
  public and deterministic --- accepting the censorship exposure of a
  predictable next leader in exchange for needing no randomness beacon.
\item \emph{Liveness recovery.} There is no view change and no
  fork-choice repair. If a matured window ever falls below $\quorum$
  blocks, the deficit can never be back-filled and the chain halts
  permanently, by design; the liveness theorems are conditional on the
  budget that has then failed, while safety is untouched. Ground this:
  at $n = 7$, $\tau \approx 4$\,s, three seats missing their slots within
  one ${\approx}\,28$\,s window halts the deployment forever.
\item \emph{Asynchronous safety.} Safety itself is clock-scoped: the
  recency rule is a clock check, so Assumption~\ref{ass:clock} is
  load-bearing for safety, not only liveness. What is never assumed is
  network synchrony --- delivery is fully adversarial in every safety
  theorem.
\end{itemize}
The trade is plain: a closed, equally weighted set with no recovery is
what lets finality be a validity predicate decided at depth $n$, rather
than a fork-choice metric resolved probabilistically over a longer
timescale.

\paragraph{Named seams.} Beyond the trusted base of
Section~\ref{sec:impl} --- the Lean kernel; the pinned translation
front-ends (Charon + Aeneas, Thales), including the hand-audited
emitter-bug corrections in the vendored TS emission; and the per-gadget
faithfulness of the circuit backend, stated rather than proved (the
equivalence lemma \code{Rust.valid\_chain\_be\_validChain} is proved at the
native u64 instantiation) --- everything not machine-checked is a named
hypothesis standing next to the theorem that consumes it: the operational
assumptions of Section~\ref{sec:assumptions}; the signature
surfaces of mode 3 and of the untimed mode-2 form (assumed, not derived from the timed model;
Assumption~\ref{ass:rotation-honest} records their operational content
directly); version independence
of keys under theft (instantiable with key-insulated
signatures~\cite{keyinsulatedsig}); the circuit path's out-of-circuit trust
anchor (the roster commitment to the cold roots); and the recency check
itself, which lives in the unverified node loop since the pure validators
carry no clock --- threading $\mathit{now}$ through them is mechanical
future work; and the per-mode operational assumptions of
Section~\ref{sec:rotation} (lockstep coordination and erasure in mode 3;
in mode 2, key-provisioning discipline is charged to the budget, not
assumed).

\paragraph{The long-range residual, and options against it.} Exposure
after a window's stretch of real time does not count against it, so a
patient adversary can bank signatures from keys exposed late --- stolen,
or reached through a signer it controls for a while --- into an
arbitrarily long fork on old stamps. Such a fork is forever stale, hence
harmless to any client applying the recency rule of
Theorem~\ref{thm:lc} (and Theorem~\ref{thm:forge} forbids preparing one in
advance), but never impossible. Two options would remove the residual
rather than contain it, each at a price: forward-secure
signatures~\cite{forwardsec} with erasure (a stateful signer); or counting
every exposure forever (\code{exposure\_agreement\_ever\_on} covers that
budget too), which charges a cleaned-up signer to the budget for good.
Key rotation does not remove it: in mode 2 a stolen current key meets the
floor of every older stamp, and in mode 3 a fork may switch to the
current generation at its branch point, so in both a late theft can still
sign old stamps --- into a fork that recency rejects. Consensus-managed
one-time keys (the XMSS variant of Section~\ref{sec:impl}) shift key state
into the chain and are complementary. We keep the windowed budget, the weakest of these
assumptions, and accept the contained residual.

\paragraph{Artifact.} The repository builds with \code{lake build}:
the model and results, both implementation bridges, and the paper-facing
\code{Molt} library with its axiom guards. A fresh checkout re-verifies
every theorem cited in this paper with that one command. That the
vendored translations are what the pinned front-ends (Charon + Aeneas;
Thales) emit from \code{rust/src/lib.rs} and \code{moltPetit.ts} is
re-checked separately by the repository's \code{nix flake check}, which
asserts faithfulness only and says nothing about the proofs. What no
machine re-checks is exactly the trusted base above: the translators and
circuit gadgets themselves, and the hand-audited re-emission
corrections.

\begin{thebibliography}{99}

\bibitem{ouroboros}
A.~Kiayias, A.~Russell, B.~David, R.~Oliynykov.
\emph{Ouroboros: A Provably Secure Proof-of-Stake Blockchain Protocol}.
CRYPTO 2017.

\bibitem{praos}
B.~David, P.~Gazi, A.~Kiayias, A.~Russell.
\emph{Ouroboros Praos: An Adaptively-Secure, Semi-synchronous
Proof-of-Stake Blockchain}. EUROCRYPT 2018.

\bibitem{samasika}
J.~Bonneau, I.~Meckler, V.~Rao, E.~Shapiro.
\emph{Ouroboros Samasika: A Proof-of-Stake Consensus Protocol for
Succinct Blockchains}. O(1)~Labs / Mina Protocol technical paper, 2020.

\bibitem{genesis}
C.~Badertscher, P.~Gazi, A.~Kiayias, A.~Russell, V.~Zikas.
\emph{Ouroboros Genesis: Composable Proof-of-Stake Blockchains with
Dynamic Availability}. ACM CCS 2018.

\bibitem{weaksubj}
V.~Buterin. \emph{Proof of Stake: How I Learned to Love Weak
Subjectivity}. Ethereum blog, 2014.

\bibitem{mina}
J.~Bonneau, I.~Meckler, V.~Rao, E.~Shapiro.
\emph{Mina: Decentralized Cryptocurrency at Scale}. Whitepaper, 2020.
IACR ePrint 2020/352.

\bibitem{plumo}
P.~Vesely, K.~G\"urkan, M.~Straka, A.~Gabizon, P.~Jovanovic,
G.~Konstantopoulos, A.~Oines, M.~Olszewski, E.~Tromer.
\emph{Plumo: An Ultralight Blockchain Client}. Financial Cryptography
2022.

\bibitem{nipopow}
A.~Kiayias, A.~Miller, D.~Zindros.
\emph{Non-Interactive Proofs of Proof-of-Work}. Financial Cryptography
2020.

\bibitem{flyclient}
B.~B\"unz, L.~Kiffer, L.~Luu, M.~Zamani.
\emph{FlyClient: Super-Light Clients for Cryptocurrencies}. IEEE S\&P
2020.

\bibitem{tendermint}
E.~Buchman, J.~Kwon, Z.~Milosevic.
\emph{The latest gossip on BFT consensus}. arXiv:1807.04938, 2018.

\bibitem{hotstuff}
M.~Yin, D.~Malkhi, M.~K.~Reiter, G.~Golan~Gueta, I.~Abraham.
\emph{HotStuff: BFT Consensus in the Lens of Blockchain}. PODC 2019.

\bibitem{canettisig}
R.~Canetti.
\emph{Universally Composable Signature, Certification, and
Authentication}. CSFW 2004.

\bibitem{forwardsec}
M.~Bellare, S.~K.~Miner.
\emph{A Forward-Secure Digital Signature Scheme}. CRYPTO 1999.

\bibitem{keyinsulated}
Y.~Dodis, J.~Katz, S.~Xu, M.~Yung.
\emph{Key-Insulated Public Key Cryptosystems}. EUROCRYPT 2002.

\bibitem{keyinsulatedsig}
Y.~Dodis, J.~Katz, S.~Xu, M.~Yung.
\emph{Strong Key-Insulated Signature Schemes}. PKC 2003.

\bibitem{stakebleeding}
P.~Ga\v{z}i, A.~Kiayias, A.~Russell.
\emph{Stake-Bleeding Attacks on Proof-of-Stake Blockchains}. Crypto Valley
Conference on Blockchain Technology (CVCBT) 2018.

\bibitem{sphincsplus}
D.~J.~Bernstein, A.~H\"ulsing, S.~K\"olbl, R.~Niederhagen, J.~Rijneveld,
P.~Schwabe.
\emph{The SPHINCS+ Signature Framework}. ACM CCS 2019.

\bibitem{fips205}
National Institute of Standards and Technology.
\emph{Stateless Hash-Based Digital Signature Standard}. FIPS 205, 2024.

\bibitem{xmssrfc}
A.~H\"ulsing, D.~Butin, S.~Gazdag, J.~Rijneveld, A.~Mohaisen.
\emph{XMSS: eXtended Merkle Signature Scheme}. RFC 8391, 2018.

\bibitem{poseidon}
L.~Grassi, D.~Khovratovich, C.~Rechberger, A.~Roy, M.~Schofnegger.
\emph{Poseidon: A New Hash Function for Zero-Knowledge Proof Systems}.
USENIX Security 2021.

\bibitem{velisarios}
V.~Rahli, I.~Vukoti\'c, M.~V\"olp, P.~Esteves-Verissimo.
\emph{Velisarios: Byzantine Fault-Tolerant Protocols Powered by Coq}.
ESOP 2018.

\bibitem{verdi}
J.~R.~Wilcox, D.~Woos, P.~Panchekha, Z.~Tatlock, X.~Wang, M.~D.~Ernst,
T.~Anderson.
\emph{Verdi: A Framework for Implementing and Formally Verifying
Distributed Systems}. PLDI 2015.

\bibitem{ironfleet}
C.~Hawblitzel, J.~Howell, M.~Kapritsos, J.~R.~Lorch, B.~Parno,
M.~L.~Roberts, S.~Setty, B.~Zill.
\emph{IronFleet: Proving Practical Distributed Systems Correct}. SOSP
2015.

\bibitem{lean}
L.~de~Moura, S.~Ullrich.
\emph{The Lean 4 Theorem Prover and Programming Language}. CADE 2021.

\bibitem{mathlib}
The mathlib Community.
\emph{The Lean Mathematical Library}. CPP 2020.

\bibitem{thales}
J.~Alama. \emph{Thales: a TypeScript-to-Lean compiler}.
\url{https://github.com/jessealama/thales}.

\bibitem{aeneas}
S.~Ho, J.~Protzenko. \emph{Aeneas: Rust Verification by Functional
Translation}. ICFP 2022. (With the Charon frontend,
\url{https://github.com/AeneasVerif}.)

\bibitem{plonky2}
Polygon Zero. \emph{Plonky2: Fast Recursive Arguments with PLONK and
FRI}. \url{https://github.com/0xPolygonZero/plonky2}.

\end{thebibliography}

\end{document}
